% !TEX program = xelatex
% 中文分析报告模板（ctexart + XeLaTeX）
% 编译：latexmk -xelatex report.tex
%
% 已经处理好的坑（不要改回去）：
%   · \setCJKmonofont  —— 否则 \texttt 里的中文会缺字（lmmono 无 CJK）
%   · 图/表用 center + \captionof（非浮动）—— 避免大图被挤到文末
%   · \raggedbottom —— 避免非浮动排版产生的巨大段间距
%   · 不要用 \path{} 包中文路径 —— hyperref 会绕过 CJK 字体映射
%   · 文件名里的下划线要写成 \_
\documentclass[UTF8, a4paper, 11pt]{ctexart}

\usepackage[margin=2.2cm]{geometry}
\usepackage{graphicx}
\usepackage{booktabs}
\usepackage{longtable}
\usepackage{array}
\usepackage{caption}
\usepackage{xcolor}
\usepackage{amsmath}
\usepackage[hidelinks]{hyperref}

\setCJKmonofont{Songti SC}     % macOS；Windows 用 SimSun，Linux 用 Noto Serif CJK SC
\setmainfont{Arial}
\setsansfont{Arial}

\captionsetup{font=small, labelfont=bf, skip=4pt}
\graphicspath{{../results/figures/pub/}}
\setlength{\parskip}{0.35em}
\renewcommand{\arraystretch}{1.15}

\title{\vspace{-1.6cm}\bfseries {{TITLE}}\\[2pt]
\large {{SUBTITLE}}}
\author{}
\date{ {{DATE}} }

\raggedbottom
\begin{document}
\maketitle
\vspace{-0.9cm}

\begin{abstract}
\noindent
{{摘要：一句话结论 + 每个突变体组的排序 + 关键倍数 + 仪器一致性。}}
\end{abstract}

\section{实验设计}

\subsection{孔位布局}

{{说明每组几列、每列几个重复、对照在哪里。}}

\begin{center}
\captionof{table}{转染板孔位布局。}
\label{tab:layout}
\small
\begin{tabular}{@{}llll@{}}
\toprule
行 & 组 1 & 组 2 & 单质粒对照\\
\midrule
B & （阴性对照） & （阴性对照） &  \\
C &  &  &  \\
D &  &  &  \\
E &  &  &  \\
F &  &  &  \\
\bottomrule
\end{tabular}
\end{center}

\subsection{数据来源}

\begin{itemize}
  \item \textbf{EVOS 10$\times$}：\texttt{/path/to/EVOS}
  \item \textbf{Spark 2$\times$}：\texttt{/path/to/Spark}
\end{itemize}

原始数据只读；分析与出图脚本见 \texttt{scripts/}。

\section{方法}

\subsection{EVOS 图像定量}

全部孔位的曝光与增益一致时可直接比较；否则按曝光$\times$增益归一化。
照明背景用大窗口中值滤波扣除。每孔计算
\begin{equation}
  \text{reporter index}=
  \frac{\sum_{\text{report mask}}\left(I_{\text{rep}}-B_{\text{rep}}\right)}
       {\sum_{\text{internal mask}}\left(I_{\text{int}}-B_{\text{int}}\right)} .
\end{equation}

\subsection{Spark 图像定量}

使用 \texttt{ProcessingPlugin/Images/*\_Processed.tiff}（已做串色矫正与背景扣除），
以空孔分布的高分位作为阳性阈值，只对超阈值像素积分：
\begin{equation}
  \text{index}=\frac{\sum_{I_{\text{rep}}>T_{\text{rep}}} I_{\text{rep}}}
                     {\sum_{I_{\text{int}}>T_{\text{int}}} I_{\text{int}}} .
\end{equation}

\subsection{相对活性与统计}

\begin{equation}
  \text{relative activity (\%)}=
  \frac{\bar{x}-\bar{x}_{\text{ctrl}}}{\bar{x}_{\text{best}}-\bar{x}_{\text{ctrl}}}\times100 .
\end{equation}
每组 $n$ 个复孔；单因素 ANOVA + Welch $t$ 检验（vs 同组对照）。

\section{结果}

\subsection{组 1}

{{文字。}}

\begin{center}
\includegraphics[width=\textwidth]{{{GROUP1_SCREEN_PDF}}}
\captionof{figure}{{\textbf{组 1 活性筛查。} a/b/c，三个定量指标；
d，代表性视野（上：报告基因；下：内参）。柱为平均值，误差线为标准差，
$n$ 个复孔；$^{*}P<0.05$，$^{**}P<0.01$（Welch $t$）。}}
\label{fig:g1}
\end{center}

\input{summary_tables.tex}

\subsection{仪器一致性}

{{Spearman ρ 与判据。}}

\begin{center}
\includegraphics[width=0.86\textwidth]{validation.pdf}
\captionof{figure}{{\textbf{仪器一致性与方法可信度。}}}
\label{fig:valid}
\end{center}

\section{讨论与局限}

\begin{enumerate}
  \item {{动态范围、重复性}}
  \item {{是否有暗对照；能否证明光依赖}}
  \item {{通道串色核验}}
  \item {{对照渗漏}}
  \item {{异常突变体需复核}}
\end{enumerate}

\section{数据与代码}

脚本 \texttt{01}–\texttt{05} 依序运行即可复现；逐孔数据见
\texttt{results/*.csv}。

\appendix
\section{逐孔原始数据}

\input{well_table.tex}

\end{document}
